%% ---------------------------------------------------------------------------
%%  lampiran.tex -- halaman lampiran, disisipkan di bagian paling bawah naskah.
%%
%%  Dipanggil dari main.tex dengan:   %% ---------------------------------------------------------------------------
%%  lampiran.tex -- halaman lampiran, disisipkan di bagian paling bawah naskah.
%%
%%  Dipanggil dari main.tex dengan:   %% ---------------------------------------------------------------------------
%%  lampiran.tex -- halaman lampiran, disisipkan di bagian paling bawah naskah.
%%
%%  Dipanggil dari main.tex dengan:   %% ---------------------------------------------------------------------------
%%  lampiran.tex -- halaman lampiran, disisipkan di bagian paling bawah naskah.
%%
%%  Dipanggil dari main.tex dengan:   \input{lampiran}
%%  tepat sebelum \end{document}, sesudah daftar pustaka.
%%
%%  >>> YANG PERLU ANDA ISI ADA DI BLOK "TAUTAN" DI BAWAH INI. <<<
%%      Ganti isi tiga makro berikut, lalu kompilasi ulang. Tidak ada tempat
%%      lain yang perlu disentuh -- seluruh tabel di bawah membaca makro ini.
%% ---------------------------------------------------------------------------

% ===== TAUTAN — ISI DI SINI =================================================

% Repositori GitHub (sudah terisi; ubah kalau repo dipindahkan).
\newcommand{\tautanRepo}{https://github.com/stephenprasetyachrismawan/openmrs-tfidf-ngram}

% Google Drive — kumpulan PPT progress. Ganti teks di dalam kurung kurawal.
\newcommand{\tautanPPT}{\textit{[tempelkan tautan Google Drive di sini]}}

% Google Drive — demo aplikasi (rekaman layar / berkas demo).
\newcommand{\tautanDemo}{\textit{[tempelkan tautan Google Drive di sini]}}

% ===== AKHIR BLOK TAUTAN ====================================================



% Judul bagian lampiran. Tidak memakai \section* bawaan ieeecolor.cls
% karena spasi vertikalnya dirancang untuk tata letak dua kolom dan
% meninggalkan lubang besar pada halaman satu kolom.
\newcommand{\lampiransec}[2]{%
  \par\addvspace{3ex}%
  \noindent{\large\bfseries Lampiran #1\quad #2}%
  \par\addvspace{1.2ex}%
  \noindent\ignorespaces}

\newpage
\onecolumn
\raggedbottom   % kelas IEEE memakai \flushbottom;
                % pada halaman lampiran yang isinya sedikit, itu meregangkan
                % jarak antar-paragraf sampai memenuhi tinggi halaman.

\begin{center}
{\Large\bfseries LAMPIRAN}\\[2pt]
{\normalsize Pencarian Terpadu Tahan Salah Ketik pada OpenMRS}
\end{center}

\vspace{1.5ex}

\noindent
Lampiran ini memuat seluruh tautan ke berkas pendukung penelitian: kode sumber,
data, tutorial menjalankan, panduan pengguna, dan bahan presentasi. Seluruh
berkas yang dapat dipublikasikan berada pada satu repositori GitHub; berkas
yang berukuran besar atau tidak cocok disimpan di repositori kode (rekaman
demo, kumpulan slide) berada pada Google Drive dan tautannya dicantumkan di
Lampiran~B.

% ---------------------------------------------------------------------------
\lampiransec{A}{Repositori Kode Sumber dan Berkas Pendukung}

\noindent
\textbf{Repositori utama:} \url{\tautanRepo}

\vspace{1ex}
\noindent
Seluruh butir pada tabel berikut berada di dalam repositori tersebut. Kolom
``Lokasi'' adalah lintasan relatif terhadap akar repositori.

\vspace{1ex}
\renewcommand{\arraystretch}{1.25}
\begin{longtable}{@{}p{0.28\textwidth}p{0.27\textwidth}p{0.41\textwidth}@{}}
\toprule
\textbf{Butir lampiran} & \textbf{Lokasi di repositori} & \textbf{Keterangan} \\
\midrule
\endfirsthead
\multicolumn{3}{@{}l}{\footnotesize\itshape Lampiran A, lanjutan}\\
\toprule
\textbf{Butir lampiran} & \textbf{Lokasi di repositori} & \textbf{Keterangan} \\
\midrule
\endhead
\midrule
\multicolumn{3}{r@{}}{\footnotesize\itshape bersambung ke halaman berikutnya}\\
\endfoot
\bottomrule
\endlastfoot

\textbf{Kode sumber modul OpenMRS} &
\texttt{backend/} \newline
\texttt{~~openmrs-module-} \newline
\texttt{~~tfidf-search/} &
Modul Java (\texttt{.omod}). Berisi seluruh algoritma K1--K6: normalisasi,
tokenisasi kata dan kepingan karakter, pembobotan \emph{ltc}, fusi tingkat
satu, \emph{Weighted RRF}, baseline B0, jalur saran ketik, serta dua endpoint
REST. Dilengkapi 47 uji unit. \\

\textbf{Kode sumber antarmuka} &
\texttt{frontend/} \newline
\texttt{~~esm-unified-search/} &
Microfrontend (ESM) untuk OpenMRS Reference Application~3: halaman
``Pencarian Terpadu'', halaman ``Perbandingan Pencarian'', dan panel
evaluasi. \\

\textbf{Kode eksperimen} &
\texttt{riset/} &
Implementasi rujukan Python untuk seluruh angka pada naskah:
\texttt{eksperimen2.py} (sistem B0--E4, himpunan uji),
\texttt{eksperimen3\_baseline\_asli.py} (baseline B0$'$),
\texttt{eksperimen\_k2.py} (jalur saran ketik),
\texttt{sapuan\_dev.py} dan \texttt{sapuan\_alpha\_dev.py} (sapuan parameter). \\

\textbf{Dataset / basis data} &
\texttt{riset/data/*.jsonl} \newline
\texttt{riset/ekspor\_*.sql} &
Korpus hasil ekspor dari basis data OpenMRS dalam format JSON Lines
(\texttt{konsep}, \texttt{obat}, \texttt{pasien}, \texttt{lain},
\texttt{hasillab}, \texttt{kondisi}), beserta kueri SQL yang dipakai untuk
mengekspornya sehingga dapat dibangkitkan ulang dari instalasi OpenMRS mana
pun. \\

\textbf{Berkas hasil eksperimen} &
\texttt{riset/hasil4/} \dots \texttt{riset/hasil6/} &
Keluaran mentah tiap eksperimen (\texttt{hasil.json},
\texttt{per\_query.json}, \texttt{ringkasan.csv}, \texttt{laporan.md}).
Pemetaan berkas ke tabel dan gambar naskah ada pada Lampiran~F. \\

\textbf{Tutorial menjalankan kode} &
\texttt{README.md} &
Panduan lengkap dari nol: prasyarat, penyalaan stack Docker, pembangunan
modul, pemasangan, verifikasi, sampai halaman terbuka di peramban. Disertai
bagian penanganan galat yang sering muncul. \\

\textbf{Panduan reproduksi hasil} &
\texttt{docs/reproduksi.md} &
Sembilan langkah untuk mengulang seluruh angka penelitian, masing-masing
dengan kriteria ``cara tahu berhasil'' yang dapat diperiksa. \\

\textbf{Panduan pengguna} &
\texttt{docs/panduan-pengguna.md} &
Cara memakai halaman Pencarian Terpadu, halaman Perbandingan Pencarian,
kotak saran ketik, dan panel evaluasi dari sudut pandang pengguna akhir. \\

\textbf{Spesifikasi algoritma} &
\texttt{docs/algoritma.md} &
Spesifikasi K1--K6, baseline B0, dan jalur saran ketik K2, cukup rinci untuk
diimplementasikan ulang tanpa menebak. \\

\textbf{Kontrak data} &
\texttt{docs/kontrak-data.md} &
Bentuk dokumen virtual yang menyatukan delapan entitas OpenMRS; wajib diikuti
persis bila menambahkan entitas baru. \\

\textbf{Catatan keputusan} &
\texttt{docs/keputusan.md} &
Seluruh keputusan rancangan dan temuan penelitian secara kronologis, termasuk
kekeliruan yang ditemukan dan cara memperbaikinya. Rujukan utama untuk
pengembangan lanjutan. \\

\textbf{Fakta lingkungan} &
\texttt{docs/lingkungan.md} &
Versi platform terverifikasi, alamat, kredensial basis data, jumlah baris
tabel, dan jebakan lingkungan yang sudah terdokumentasi. \\

\textbf{Aturan proyek} &
\texttt{CLAUDE.md} &
Sepuluh aturan yang mengikat seluruh perubahan kode dan dokumen, termasuk
aturan determinisme dan larangan mengubah angka hasil. \\

\textbf{Proposal penelitian} &
\texttt{docs/proposal.html} &
Proposal lengkap beserta \emph{mockup} interaktif dan contoh perhitungan. \\

\textbf{Naskah artikel ini} &
\texttt{article/} &
Sumber \LaTeX{} naskah, skrip pembangun gambar, dan berkas hasil kompilasi. \\

\end{longtable}

% ---------------------------------------------------------------------------
\vspace{3ex}
\lampiransec{B}{Tautan Berkas Pendukung Eksternal}

\noindent
Berkas berikut berada di Google Drive karena ukurannya tidak sesuai untuk
disimpan pada repositori kode.

\vspace{1ex}
\renewcommand{\arraystretch}{1.6}
\begin{tabular}{@{}p{0.24\textwidth}p{0.72\textwidth}@{}}
\toprule
\textbf{Berkas} & \textbf{Tautan} \\
\midrule

\textbf{Kumpulan PPT Progress} \newline
{\footnotesize Seluruh slide laporan kemajuan penelitian, berurutan menurut
tanggal.} &
\tautanPPT \\

\textbf{Demo Aplikasi} \newline
{\footnotesize Rekaman layar modul yang berjalan: kueri gagal pada kotak
pencarian bawaan, kueri yang sama berhasil pada sistem usulan.} &
\tautanDemo \\

\bottomrule
\end{tabular}

\vspace{1.5ex}
\noindent
{\footnotesize \textbf{Catatan bagi penulis:} kedua tautan di atas diisi pada
blok \texttt{TAUTAN} di kepala berkas \texttt{article/lampiran.tex}
(makro \texttt{\textbackslash tautanPPT} dan \texttt{\textbackslash
tautanDemo}). Tidak perlu mengubah tabel ini.}

% ---------------------------------------------------------------------------
\vspace{3ex}
\lampiransec{C}{Peta Isi Repositori}

\begin{verbatim}
openmrs-tfidf-ngram/
|
+-- README.md                 tutorial lengkap menjalankan dari nol
+-- CLAUDE.md                 aturan proyek yang mengikat
|
+-- backend/openmrs-module-tfidf-search/
|   +-- api/                  algoritma inti + 47 uji unit
|   |   +-- src/main/java/org/openmrs/module/unifiedsearch/
|   |   |   +-- TextNormalizer.java        normalisasi teks
|   |   |   +-- Tokenizer.java             token kata + kepingan karakter
|   |   |   +-- TfIdfIndex.java            indeks terbalik, pembobotan ltc
|   |   |   +-- FusionSearch.java          fusi tingkat 1 (alpha)
|   |   |   +-- WeightedRrf.java           fusi tingkat 2 (Weighted RRF)
|   |   |   +-- OpenMrsHeuristic.java      baseline B0
|   |   |   +-- BigramJaccardSuggester.java  jalur saran ketik
|   |   |   +-- RankingEngine.java         pemilih mode b0/b1/e1/e3
|   |   |   +-- IndexBuilder.java          pembangunan 13 indeks saat startup
|   |   |   +-- EvalService.java           endpoint evaluasi
|   |   |   +-- source/                    proyeksi 8 tabel OpenMRS
|   |   +-- src/main/resources/gold-dev-100.json   standar emas 100 kueri dev
|   +-- omod/                 lapisan web: REST controller, halaman JSP
|
+-- frontend/esm-unified-search/   microfrontend RefApp 3
|
+-- riset/                    eksperimen Python
|   +-- eksperimen2.py                sistem B0-E4, himpunan uji 180 kueri
|   +-- eksperimen3_baseline_asli.py  baseline B0' (OpenMRS asli)
|   +-- eksperimen_k2.py              jalur saran ketik
|   +-- sapuan_dev.py                 sapuan NGRAM / K_RRF / EPS
|   +-- sapuan_alpha_dev.py           sapuan ALPHA
|   +-- data/*.jsonl                  korpus hasil ekspor
|   +-- ekspor_*.sql                  kueri SQL pembangkit korpus
|   +-- hasil4..hasil6/               keluaran mentah tiap eksperimen
|
+-- docs/                     dokumentasi
|   +-- reproduksi.md         9 langkah mengulang hasil
|   +-- panduan-pengguna.md   panduan pengguna akhir
|   +-- algoritma.md          spesifikasi K1-K6 + B0 + K2
|   +-- kontrak-data.md       bentuk dokumen virtual
|   +-- keputusan.md          catatan keputusan dan temuan, kronologis
|   +-- lingkungan.md         versi, alamat, kredensial, jebakan
|   +-- proposal.html         proposal + mockup interaktif
|
+-- scripts/                  otomasi pemasangan (PowerShell)
|   +-- pasang-modul.ps1      build .omod -> salin -> restart -> tunggu sehat
|   +-- pasang-esm.ps1        build ESM -> sisipkan ke importmap
|   +-- verify-openmrs.ps1    pemeriksaan kesehatan menyeluruh
|
+-- article/                  naskah artikel ini
|   +-- main.tex, lampiran.tex, referensi.bib
|   +-- gambar/buat_gambar.py  membangun gambar dari berkas hasil
|
+-- tugas/                    14 berkas tugas berurutan, masing-masing
                              dengan kriteria uji "Selesai kalau"
\end{verbatim}

% ---------------------------------------------------------------------------
\newpage
\lampiransec{D}{Ringkasan Langkah Menjalankan}

\noindent
Ringkasan ini memuat perintah intinya saja. Penjelasan lengkap setiap langkah,
termasuk cara memastikan langkah itu berhasil dan penanganan galat yang sering
muncul, ada pada \texttt{README.md} di akar repositori.

\vspace{1ex}
\noindent\textbf{Prasyarat:} Docker Desktop, Git, JDK~17, Node.js~20+,
Python~3.9+, dan PowerShell (Windows). Rincian versi terverifikasi ada pada
\texttt{docs/lingkungan.md}.

\vspace{1.5ex}
\renewcommand{\arraystretch}{1.2}
\begin{tabular}{@{}p{0.05\textwidth}p{0.40\textwidth}p{0.51\textwidth}@{}}
\toprule
\textbf{No.} & \textbf{Langkah} & \textbf{Perintah} \\
\midrule
1 & Klon repositori penelitian &
\texttt{git clone https://github.com/} \newline
\texttt{~~stephenprasetyachrismawan/} \newline
\texttt{~~openmrs-tfidf-ngram.git} \\
2 & Klon stack Docker OpenMRS sejajar dengan repositori penelitian &
\texttt{git clone https://github.com/} \newline
\texttt{openmrs/openmrs-distro-} \newline
\texttt{referenceapplication.git} \\
3 & Nyalakan stack (tanpa \texttt{-p}) &
\texttt{cd openmrs-distro-referenceapplication} \newline
\texttt{docker compose up -d} \\
4 & Tunggu empat kontainer sehat &
\texttt{docker ps} \\
5 & Bangun dan pasang modul backend &
\texttt{.\textbackslash scripts\textbackslash pasang-modul.ps1} \\
6 & Bangun dan pasang antarmuka ESM &
\texttt{.\textbackslash scripts\textbackslash pasang-esm.ps1} \\
7 & Verifikasi menyeluruh &
\texttt{.\textbackslash scripts\textbackslash verify-openmrs.ps1} \\
8 & Buka di peramban &
\texttt{http://127.0.0.1/openmrs/spa/} \newline
\texttt{unified-search} \\
\bottomrule
\end{tabular}

\vspace{2ex}
\noindent
\textbf{Peringatan lingkungan.} Seluruh alamat memakai \texttt{http://127.0.0.1},
\emph{bukan} \texttt{localhost}. Pada sebagian mesin Windows, \texttt{localhost}
teresolusi ke IPv6 \texttt{::1} yang mungkin dikuasai server lain, sehingga
seluruh lintasan \texttt{/openmrs/\dots} mengembalikan 404. Bila muncul halaman
bawaan Apache, yang salah adalah alamatnya, bukan modulnya.

\vspace{1ex}
\noindent
\textbf{Kredensial demo bawaan:} \texttt{admin} / \texttt{Admin123}.

% ---------------------------------------------------------------------------
\vspace{3ex}
\lampiransec{E}{Antarmuka REST Modul}

\noindent
Basis alamat: \texttt{http://127.0.0.1/openmrs/ws/rest/v1/unifiedsearch}

\vspace{1.5ex}
\renewcommand{\arraystretch}{1.3}
\begin{tabular}{@{}p{0.20\textwidth}p{0.34\textwidth}p{0.42\textwidth}@{}}
\toprule
\textbf{Endpoint} & \textbf{Parameter} & \textbf{Keterangan} \\
\midrule

\texttt{GET /unifiedsearch} &
\texttt{q} (wajib) \newline
\texttt{mode} = \texttt{b0} $|$ \texttt{b1} $|$ \texttt{e1} $|$ \texttt{e3}
  (bawaan \texttt{e3}) \newline
\texttt{limit} (bawaan 10) \newline
\texttt{entitas} (opsional) &
Pencarian terpadu. Parameter \texttt{mode} memilih sistem, sehingga kontribusi
tiap komponen dapat diukur ulang kapan saja tanpa mengubah kode. \\

\texttt{GET /unifiedsearch/} \newline \texttt{saran} &
\texttt{q} (wajib) \newline
\texttt{limit} (bawaan 6) &
Saran ketik ``Maksud Anda'' berbasis Jaccard bigram. Jalur terpisah; tidak
menyentuh peringkat. Hasil berjenis pasien, hasil laboratorium, dan kondisi
disaring menurut privilese \texttt{View Patients} pemanggil. \\

\texttt{GET /unifiedsearch/} \newline \texttt{eval} &
\texttt{mode} (bawaan \texttt{e3}) &
Menjalankan evaluasi atas 100 kueri pengembangan dan mengembalikan P@1, P@5,
R@10, MRR, MAP, nDCG@10, serta proporsi kueri tanpa hasil. Jawaban menyertakan
\texttt{gold\_sha256} agar standar emas yang dipakai dapat dicocokkan. \\

\bottomrule
\end{tabular}

\vspace{1.5ex}
\noindent
Setiap jawaban menyertakan tajuk \texttt{X-Unifiedsearch-Waktu-Ms} berisi waktu
pemrosesan. Waktu sengaja ditempatkan pada tajuk, bukan pada badan JSON, agar
dua panggilan yang sama menghasilkan badan jawaban yang identik pada tingkat
bita --- syarat pengujian determinisme.

\vspace{1.5ex}
\noindent\textbf{Contoh pemanggilan:}
\begin{verbatim}
curl -u admin:Admin123 \
  "http://127.0.0.1/openmrs/ws/rest/v1/unifiedsearch?q=diabete%20melitus&mode=e3"
\end{verbatim}

% ---------------------------------------------------------------------------
\vspace{2ex}
\lampiransec{F}{Pemetaan Berkas Hasil ke Tabel dan Gambar Naskah}

\noindent
Tidak ada angka pada naskah ini yang diketik ulang secara manual ke dalam
gambar. Skrip \texttt{article/gambar/buat\_gambar.py} membaca berkas hasil
di bawah dan membangun ulang seluruh gambar; bila eksperimen dijalankan ulang
dan angkanya berubah, gambarnya ikut berubah.

\vspace{1.5ex}
\renewcommand{\arraystretch}{1.25}
\begin{tabular}{@{}p{0.26\textwidth}p{0.30\textwidth}p{0.40\textwidth}@{}}
\toprule
\textbf{Bagian naskah} & \textbf{Berkas sumber} & \textbf{Cakupan} \\
\midrule
Tabel~\ref{tab:utama}, \ref{tab:ablasi}, \ref{tab:tipe}, \ref{tab:entitas};
Gambar~\ref{fig:sistem}, \ref{fig:tipe} &
\texttt{riset/hasil6/hasil.json} &
260 kueri uji, korpus 8 entitas, sekali jalan, parameter terkunci \\

Gambar~\ref{fig:sapuan} &
\texttt{riset/hasil6/sapuan\_dev.json} &
Sapuan $\alpha$, $n$, $k$, $\varepsilon$ pada 100 kueri pengembangan \\

Tabel~\ref{tab:b0prime}; Gambar~\ref{fig:b0prime} &
\texttt{riset/hasil4/hasil.json} &
34 kueri konsep, baseline B0$'$ (endpoint OpenMRS asli) \\

Tabel~\ref{tab:k2}; Gambar~\ref{fig:k2} &
\texttt{riset/hasil5/hasil.json} &
214 kueri, jalur saran ketik, korpus 8 entitas \\

Tabel~\ref{tab:verifikasi}, angka latensi dan determinisme &
\texttt{docs/keputusan.md} bagian C1--C4, F1--F2 &
Verifikasi silang Java--Python dan pengukuran kinerja modul \\
\bottomrule
\end{tabular}

\vspace{2ex}
\noindent
\textbf{Cara membangkitkan ulang seluruh angka:} jalankan
\texttt{python riset/eksperimen2.py} (himpunan uji),
\texttt{python riset/eksperimen3\_baseline\_asli.py} (baseline B0$'$, perlu
stack Docker berjalan), dan \texttt{python riset/eksperimen\_k2.py}
(jalur saran ketik), lalu \texttt{python article/gambar/buat\_gambar.py} untuk
membangun ulang gambar. Seluruh skrip deterministik: benih acak dipatok,
pengurutan memakai kunci majemuk, dan iterasi himpunan selalu diurutkan lebih
dahulu, sehingga dua proses terpisah menghasilkan berkas yang identik pada
tingkat bita.

\twocolumn

%%  tepat sebelum \end{document}, sesudah daftar pustaka.
%%
%%  >>> YANG PERLU ANDA ISI ADA DI BLOK "TAUTAN" DI BAWAH INI. <<<
%%      Ganti isi tiga makro berikut, lalu kompilasi ulang. Tidak ada tempat
%%      lain yang perlu disentuh -- seluruh tabel di bawah membaca makro ini.
%% ---------------------------------------------------------------------------

% ===== TAUTAN — ISI DI SINI =================================================

% Repositori GitHub (sudah terisi; ubah kalau repo dipindahkan).
\newcommand{\tautanRepo}{https://github.com/stephenprasetyachrismawan/openmrs-tfidf-ngram}

% Google Drive — kumpulan PPT progress. Ganti teks di dalam kurung kurawal.
\newcommand{\tautanPPT}{\textit{[tempelkan tautan Google Drive di sini]}}

% Google Drive — demo aplikasi (rekaman layar / berkas demo).
\newcommand{\tautanDemo}{\textit{[tempelkan tautan Google Drive di sini]}}

% ===== AKHIR BLOK TAUTAN ====================================================



% Judul bagian lampiran. Tidak memakai \section* bawaan ieeecolor.cls
% karena spasi vertikalnya dirancang untuk tata letak dua kolom dan
% meninggalkan lubang besar pada halaman satu kolom.
\newcommand{\lampiransec}[2]{%
  \par\addvspace{3ex}%
  \noindent{\large\bfseries Lampiran #1\quad #2}%
  \par\addvspace{1.2ex}%
  \noindent\ignorespaces}

\newpage
\onecolumn
\raggedbottom   % kelas IEEE memakai \flushbottom;
                % pada halaman lampiran yang isinya sedikit, itu meregangkan
                % jarak antar-paragraf sampai memenuhi tinggi halaman.

\begin{center}
{\Large\bfseries LAMPIRAN}\\[2pt]
{\normalsize Pencarian Terpadu Tahan Salah Ketik pada OpenMRS}
\end{center}

\vspace{1.5ex}

\noindent
Lampiran ini memuat seluruh tautan ke berkas pendukung penelitian: kode sumber,
data, tutorial menjalankan, panduan pengguna, dan bahan presentasi. Seluruh
berkas yang dapat dipublikasikan berada pada satu repositori GitHub; berkas
yang berukuran besar atau tidak cocok disimpan di repositori kode (rekaman
demo, kumpulan slide) berada pada Google Drive dan tautannya dicantumkan di
Lampiran~B.

% ---------------------------------------------------------------------------
\lampiransec{A}{Repositori Kode Sumber dan Berkas Pendukung}

\noindent
\textbf{Repositori utama:} \url{\tautanRepo}

\vspace{1ex}
\noindent
Seluruh butir pada tabel berikut berada di dalam repositori tersebut. Kolom
``Lokasi'' adalah lintasan relatif terhadap akar repositori.

\vspace{1ex}
\renewcommand{\arraystretch}{1.25}
\begin{longtable}{@{}p{0.28\textwidth}p{0.27\textwidth}p{0.41\textwidth}@{}}
\toprule
\textbf{Butir lampiran} & \textbf{Lokasi di repositori} & \textbf{Keterangan} \\
\midrule
\endfirsthead
\multicolumn{3}{@{}l}{\footnotesize\itshape Lampiran A, lanjutan}\\
\toprule
\textbf{Butir lampiran} & \textbf{Lokasi di repositori} & \textbf{Keterangan} \\
\midrule
\endhead
\midrule
\multicolumn{3}{r@{}}{\footnotesize\itshape bersambung ke halaman berikutnya}\\
\endfoot
\bottomrule
\endlastfoot

\textbf{Kode sumber modul OpenMRS} &
\texttt{backend/} \newline
\texttt{~~openmrs-module-} \newline
\texttt{~~tfidf-search/} &
Modul Java (\texttt{.omod}). Berisi seluruh algoritma K1--K6: normalisasi,
tokenisasi kata dan kepingan karakter, pembobotan \emph{ltc}, fusi tingkat
satu, \emph{Weighted RRF}, baseline B0, jalur saran ketik, serta dua endpoint
REST. Dilengkapi 47 uji unit. \\

\textbf{Kode sumber antarmuka} &
\texttt{frontend/} \newline
\texttt{~~esm-unified-search/} &
Microfrontend (ESM) untuk OpenMRS Reference Application~3: halaman
``Pencarian Terpadu'', halaman ``Perbandingan Pencarian'', dan panel
evaluasi. \\

\textbf{Kode eksperimen} &
\texttt{riset/} &
Implementasi rujukan Python untuk seluruh angka pada naskah:
\texttt{eksperimen2.py} (sistem B0--E4, himpunan uji),
\texttt{eksperimen3\_baseline\_asli.py} (baseline B0$'$),
\texttt{eksperimen\_k2.py} (jalur saran ketik),
\texttt{sapuan\_dev.py} dan \texttt{sapuan\_alpha\_dev.py} (sapuan parameter). \\

\textbf{Dataset / basis data} &
\texttt{riset/data/*.jsonl} \newline
\texttt{riset/ekspor\_*.sql} &
Korpus hasil ekspor dari basis data OpenMRS dalam format JSON Lines
(\texttt{konsep}, \texttt{obat}, \texttt{pasien}, \texttt{lain},
\texttt{hasillab}, \texttt{kondisi}), beserta kueri SQL yang dipakai untuk
mengekspornya sehingga dapat dibangkitkan ulang dari instalasi OpenMRS mana
pun. \\

\textbf{Berkas hasil eksperimen} &
\texttt{riset/hasil4/} \dots \texttt{riset/hasil6/} &
Keluaran mentah tiap eksperimen (\texttt{hasil.json},
\texttt{per\_query.json}, \texttt{ringkasan.csv}, \texttt{laporan.md}).
Pemetaan berkas ke tabel dan gambar naskah ada pada Lampiran~F. \\

\textbf{Tutorial menjalankan kode} &
\texttt{README.md} &
Panduan lengkap dari nol: prasyarat, penyalaan stack Docker, pembangunan
modul, pemasangan, verifikasi, sampai halaman terbuka di peramban. Disertai
bagian penanganan galat yang sering muncul. \\

\textbf{Panduan reproduksi hasil} &
\texttt{docs/reproduksi.md} &
Sembilan langkah untuk mengulang seluruh angka penelitian, masing-masing
dengan kriteria ``cara tahu berhasil'' yang dapat diperiksa. \\

\textbf{Panduan pengguna} &
\texttt{docs/panduan-pengguna.md} &
Cara memakai halaman Pencarian Terpadu, halaman Perbandingan Pencarian,
kotak saran ketik, dan panel evaluasi dari sudut pandang pengguna akhir. \\

\textbf{Spesifikasi algoritma} &
\texttt{docs/algoritma.md} &
Spesifikasi K1--K6, baseline B0, dan jalur saran ketik K2, cukup rinci untuk
diimplementasikan ulang tanpa menebak. \\

\textbf{Kontrak data} &
\texttt{docs/kontrak-data.md} &
Bentuk dokumen virtual yang menyatukan delapan entitas OpenMRS; wajib diikuti
persis bila menambahkan entitas baru. \\

\textbf{Catatan keputusan} &
\texttt{docs/keputusan.md} &
Seluruh keputusan rancangan dan temuan penelitian secara kronologis, termasuk
kekeliruan yang ditemukan dan cara memperbaikinya. Rujukan utama untuk
pengembangan lanjutan. \\

\textbf{Fakta lingkungan} &
\texttt{docs/lingkungan.md} &
Versi platform terverifikasi, alamat, kredensial basis data, jumlah baris
tabel, dan jebakan lingkungan yang sudah terdokumentasi. \\

\textbf{Aturan proyek} &
\texttt{CLAUDE.md} &
Sepuluh aturan yang mengikat seluruh perubahan kode dan dokumen, termasuk
aturan determinisme dan larangan mengubah angka hasil. \\

\textbf{Proposal penelitian} &
\texttt{docs/proposal.html} &
Proposal lengkap beserta \emph{mockup} interaktif dan contoh perhitungan. \\

\textbf{Naskah artikel ini} &
\texttt{article/} &
Sumber \LaTeX{} naskah, skrip pembangun gambar, dan berkas hasil kompilasi. \\

\end{longtable}

% ---------------------------------------------------------------------------
\vspace{3ex}
\lampiransec{B}{Tautan Berkas Pendukung Eksternal}

\noindent
Berkas berikut berada di Google Drive karena ukurannya tidak sesuai untuk
disimpan pada repositori kode.

\vspace{1ex}
\renewcommand{\arraystretch}{1.6}
\begin{tabular}{@{}p{0.24\textwidth}p{0.72\textwidth}@{}}
\toprule
\textbf{Berkas} & \textbf{Tautan} \\
\midrule

\textbf{Kumpulan PPT Progress} \newline
{\footnotesize Seluruh slide laporan kemajuan penelitian, berurutan menurut
tanggal.} &
\tautanPPT \\

\textbf{Demo Aplikasi} \newline
{\footnotesize Rekaman layar modul yang berjalan: kueri gagal pada kotak
pencarian bawaan, kueri yang sama berhasil pada sistem usulan.} &
\tautanDemo \\

\bottomrule
\end{tabular}

\vspace{1.5ex}
\noindent
{\footnotesize \textbf{Catatan bagi penulis:} kedua tautan di atas diisi pada
blok \texttt{TAUTAN} di kepala berkas \texttt{article/lampiran.tex}
(makro \texttt{\textbackslash tautanPPT} dan \texttt{\textbackslash
tautanDemo}). Tidak perlu mengubah tabel ini.}

% ---------------------------------------------------------------------------
\vspace{3ex}
\lampiransec{C}{Peta Isi Repositori}

\begin{verbatim}
openmrs-tfidf-ngram/
|
+-- README.md                 tutorial lengkap menjalankan dari nol
+-- CLAUDE.md                 aturan proyek yang mengikat
|
+-- backend/openmrs-module-tfidf-search/
|   +-- api/                  algoritma inti + 47 uji unit
|   |   +-- src/main/java/org/openmrs/module/unifiedsearch/
|   |   |   +-- TextNormalizer.java        normalisasi teks
|   |   |   +-- Tokenizer.java             token kata + kepingan karakter
|   |   |   +-- TfIdfIndex.java            indeks terbalik, pembobotan ltc
|   |   |   +-- FusionSearch.java          fusi tingkat 1 (alpha)
|   |   |   +-- WeightedRrf.java           fusi tingkat 2 (Weighted RRF)
|   |   |   +-- OpenMrsHeuristic.java      baseline B0
|   |   |   +-- BigramJaccardSuggester.java  jalur saran ketik
|   |   |   +-- RankingEngine.java         pemilih mode b0/b1/e1/e3
|   |   |   +-- IndexBuilder.java          pembangunan 13 indeks saat startup
|   |   |   +-- EvalService.java           endpoint evaluasi
|   |   |   +-- source/                    proyeksi 8 tabel OpenMRS
|   |   +-- src/main/resources/gold-dev-100.json   standar emas 100 kueri dev
|   +-- omod/                 lapisan web: REST controller, halaman JSP
|
+-- frontend/esm-unified-search/   microfrontend RefApp 3
|
+-- riset/                    eksperimen Python
|   +-- eksperimen2.py                sistem B0-E4, himpunan uji 180 kueri
|   +-- eksperimen3_baseline_asli.py  baseline B0' (OpenMRS asli)
|   +-- eksperimen_k2.py              jalur saran ketik
|   +-- sapuan_dev.py                 sapuan NGRAM / K_RRF / EPS
|   +-- sapuan_alpha_dev.py           sapuan ALPHA
|   +-- data/*.jsonl                  korpus hasil ekspor
|   +-- ekspor_*.sql                  kueri SQL pembangkit korpus
|   +-- hasil4..hasil6/               keluaran mentah tiap eksperimen
|
+-- docs/                     dokumentasi
|   +-- reproduksi.md         9 langkah mengulang hasil
|   +-- panduan-pengguna.md   panduan pengguna akhir
|   +-- algoritma.md          spesifikasi K1-K6 + B0 + K2
|   +-- kontrak-data.md       bentuk dokumen virtual
|   +-- keputusan.md          catatan keputusan dan temuan, kronologis
|   +-- lingkungan.md         versi, alamat, kredensial, jebakan
|   +-- proposal.html         proposal + mockup interaktif
|
+-- scripts/                  otomasi pemasangan (PowerShell)
|   +-- pasang-modul.ps1      build .omod -> salin -> restart -> tunggu sehat
|   +-- pasang-esm.ps1        build ESM -> sisipkan ke importmap
|   +-- verify-openmrs.ps1    pemeriksaan kesehatan menyeluruh
|
+-- article/                  naskah artikel ini
|   +-- main.tex, lampiran.tex, referensi.bib
|   +-- gambar/buat_gambar.py  membangun gambar dari berkas hasil
|
+-- tugas/                    14 berkas tugas berurutan, masing-masing
                              dengan kriteria uji "Selesai kalau"
\end{verbatim}

% ---------------------------------------------------------------------------
\newpage
\lampiransec{D}{Ringkasan Langkah Menjalankan}

\noindent
Ringkasan ini memuat perintah intinya saja. Penjelasan lengkap setiap langkah,
termasuk cara memastikan langkah itu berhasil dan penanganan galat yang sering
muncul, ada pada \texttt{README.md} di akar repositori.

\vspace{1ex}
\noindent\textbf{Prasyarat:} Docker Desktop, Git, JDK~17, Node.js~20+,
Python~3.9+, dan PowerShell (Windows). Rincian versi terverifikasi ada pada
\texttt{docs/lingkungan.md}.

\vspace{1.5ex}
\renewcommand{\arraystretch}{1.2}
\begin{tabular}{@{}p{0.05\textwidth}p{0.40\textwidth}p{0.51\textwidth}@{}}
\toprule
\textbf{No.} & \textbf{Langkah} & \textbf{Perintah} \\
\midrule
1 & Klon repositori penelitian &
\texttt{git clone https://github.com/} \newline
\texttt{~~stephenprasetyachrismawan/} \newline
\texttt{~~openmrs-tfidf-ngram.git} \\
2 & Klon stack Docker OpenMRS sejajar dengan repositori penelitian &
\texttt{git clone https://github.com/} \newline
\texttt{openmrs/openmrs-distro-} \newline
\texttt{referenceapplication.git} \\
3 & Nyalakan stack (tanpa \texttt{-p}) &
\texttt{cd openmrs-distro-referenceapplication} \newline
\texttt{docker compose up -d} \\
4 & Tunggu empat kontainer sehat &
\texttt{docker ps} \\
5 & Bangun dan pasang modul backend &
\texttt{.\textbackslash scripts\textbackslash pasang-modul.ps1} \\
6 & Bangun dan pasang antarmuka ESM &
\texttt{.\textbackslash scripts\textbackslash pasang-esm.ps1} \\
7 & Verifikasi menyeluruh &
\texttt{.\textbackslash scripts\textbackslash verify-openmrs.ps1} \\
8 & Buka di peramban &
\texttt{http://127.0.0.1/openmrs/spa/} \newline
\texttt{unified-search} \\
\bottomrule
\end{tabular}

\vspace{2ex}
\noindent
\textbf{Peringatan lingkungan.} Seluruh alamat memakai \texttt{http://127.0.0.1},
\emph{bukan} \texttt{localhost}. Pada sebagian mesin Windows, \texttt{localhost}
teresolusi ke IPv6 \texttt{::1} yang mungkin dikuasai server lain, sehingga
seluruh lintasan \texttt{/openmrs/\dots} mengembalikan 404. Bila muncul halaman
bawaan Apache, yang salah adalah alamatnya, bukan modulnya.

\vspace{1ex}
\noindent
\textbf{Kredensial demo bawaan:} \texttt{admin} / \texttt{Admin123}.

% ---------------------------------------------------------------------------
\vspace{3ex}
\lampiransec{E}{Antarmuka REST Modul}

\noindent
Basis alamat: \texttt{http://127.0.0.1/openmrs/ws/rest/v1/unifiedsearch}

\vspace{1.5ex}
\renewcommand{\arraystretch}{1.3}
\begin{tabular}{@{}p{0.20\textwidth}p{0.34\textwidth}p{0.42\textwidth}@{}}
\toprule
\textbf{Endpoint} & \textbf{Parameter} & \textbf{Keterangan} \\
\midrule

\texttt{GET /unifiedsearch} &
\texttt{q} (wajib) \newline
\texttt{mode} = \texttt{b0} $|$ \texttt{b1} $|$ \texttt{e1} $|$ \texttt{e3}
  (bawaan \texttt{e3}) \newline
\texttt{limit} (bawaan 10) \newline
\texttt{entitas} (opsional) &
Pencarian terpadu. Parameter \texttt{mode} memilih sistem, sehingga kontribusi
tiap komponen dapat diukur ulang kapan saja tanpa mengubah kode. \\

\texttt{GET /unifiedsearch/} \newline \texttt{saran} &
\texttt{q} (wajib) \newline
\texttt{limit} (bawaan 6) &
Saran ketik ``Maksud Anda'' berbasis Jaccard bigram. Jalur terpisah; tidak
menyentuh peringkat. Hasil berjenis pasien, hasil laboratorium, dan kondisi
disaring menurut privilese \texttt{View Patients} pemanggil. \\

\texttt{GET /unifiedsearch/} \newline \texttt{eval} &
\texttt{mode} (bawaan \texttt{e3}) &
Menjalankan evaluasi atas 100 kueri pengembangan dan mengembalikan P@1, P@5,
R@10, MRR, MAP, nDCG@10, serta proporsi kueri tanpa hasil. Jawaban menyertakan
\texttt{gold\_sha256} agar standar emas yang dipakai dapat dicocokkan. \\

\bottomrule
\end{tabular}

\vspace{1.5ex}
\noindent
Setiap jawaban menyertakan tajuk \texttt{X-Unifiedsearch-Waktu-Ms} berisi waktu
pemrosesan. Waktu sengaja ditempatkan pada tajuk, bukan pada badan JSON, agar
dua panggilan yang sama menghasilkan badan jawaban yang identik pada tingkat
bita --- syarat pengujian determinisme.

\vspace{1.5ex}
\noindent\textbf{Contoh pemanggilan:}
\begin{verbatim}
curl -u admin:Admin123 \
  "http://127.0.0.1/openmrs/ws/rest/v1/unifiedsearch?q=diabete%20melitus&mode=e3"
\end{verbatim}

% ---------------------------------------------------------------------------
\vspace{2ex}
\lampiransec{F}{Pemetaan Berkas Hasil ke Tabel dan Gambar Naskah}

\noindent
Tidak ada angka pada naskah ini yang diketik ulang secara manual ke dalam
gambar. Skrip \texttt{article/gambar/buat\_gambar.py} membaca berkas hasil
di bawah dan membangun ulang seluruh gambar; bila eksperimen dijalankan ulang
dan angkanya berubah, gambarnya ikut berubah.

\vspace{1.5ex}
\renewcommand{\arraystretch}{1.25}
\begin{tabular}{@{}p{0.26\textwidth}p{0.30\textwidth}p{0.40\textwidth}@{}}
\toprule
\textbf{Bagian naskah} & \textbf{Berkas sumber} & \textbf{Cakupan} \\
\midrule
Tabel~\ref{tab:utama}, \ref{tab:ablasi}, \ref{tab:tipe}, \ref{tab:entitas};
Gambar~\ref{fig:sistem}, \ref{fig:tipe} &
\texttt{riset/hasil6/hasil.json} &
260 kueri uji, korpus 8 entitas, sekali jalan, parameter terkunci \\

Gambar~\ref{fig:sapuan} &
\texttt{riset/hasil6/sapuan\_dev.json} &
Sapuan $\alpha$, $n$, $k$, $\varepsilon$ pada 100 kueri pengembangan \\

Tabel~\ref{tab:b0prime}; Gambar~\ref{fig:b0prime} &
\texttt{riset/hasil4/hasil.json} &
34 kueri konsep, baseline B0$'$ (endpoint OpenMRS asli) \\

Tabel~\ref{tab:k2}; Gambar~\ref{fig:k2} &
\texttt{riset/hasil5/hasil.json} &
214 kueri, jalur saran ketik, korpus 8 entitas \\

Tabel~\ref{tab:verifikasi}, angka latensi dan determinisme &
\texttt{docs/keputusan.md} bagian C1--C4, F1--F2 &
Verifikasi silang Java--Python dan pengukuran kinerja modul \\
\bottomrule
\end{tabular}

\vspace{2ex}
\noindent
\textbf{Cara membangkitkan ulang seluruh angka:} jalankan
\texttt{python riset/eksperimen2.py} (himpunan uji),
\texttt{python riset/eksperimen3\_baseline\_asli.py} (baseline B0$'$, perlu
stack Docker berjalan), dan \texttt{python riset/eksperimen\_k2.py}
(jalur saran ketik), lalu \texttt{python article/gambar/buat\_gambar.py} untuk
membangun ulang gambar. Seluruh skrip deterministik: benih acak dipatok,
pengurutan memakai kunci majemuk, dan iterasi himpunan selalu diurutkan lebih
dahulu, sehingga dua proses terpisah menghasilkan berkas yang identik pada
tingkat bita.

\twocolumn

%%  tepat sebelum \end{document}, sesudah daftar pustaka.
%%
%%  >>> YANG PERLU ANDA ISI ADA DI BLOK "TAUTAN" DI BAWAH INI. <<<
%%      Ganti isi tiga makro berikut, lalu kompilasi ulang. Tidak ada tempat
%%      lain yang perlu disentuh -- seluruh tabel di bawah membaca makro ini.
%% ---------------------------------------------------------------------------

% ===== TAUTAN — ISI DI SINI =================================================

% Repositori GitHub (sudah terisi; ubah kalau repo dipindahkan).
\newcommand{\tautanRepo}{https://github.com/stephenprasetyachrismawan/openmrs-tfidf-ngram}

% Google Drive — kumpulan PPT progress. Ganti teks di dalam kurung kurawal.
\newcommand{\tautanPPT}{\textit{[tempelkan tautan Google Drive di sini]}}

% Google Drive — demo aplikasi (rekaman layar / berkas demo).
\newcommand{\tautanDemo}{\textit{[tempelkan tautan Google Drive di sini]}}

% ===== AKHIR BLOK TAUTAN ====================================================



% Judul bagian lampiran. Tidak memakai \section* bawaan ieeecolor.cls
% karena spasi vertikalnya dirancang untuk tata letak dua kolom dan
% meninggalkan lubang besar pada halaman satu kolom.
\newcommand{\lampiransec}[2]{%
  \par\addvspace{3ex}%
  \noindent{\large\bfseries Lampiran #1\quad #2}%
  \par\addvspace{1.2ex}%
  \noindent\ignorespaces}

\newpage
\onecolumn
\raggedbottom   % kelas IEEE memakai \flushbottom;
                % pada halaman lampiran yang isinya sedikit, itu meregangkan
                % jarak antar-paragraf sampai memenuhi tinggi halaman.

\begin{center}
{\Large\bfseries LAMPIRAN}\\[2pt]
{\normalsize Pencarian Terpadu Tahan Salah Ketik pada OpenMRS}
\end{center}

\vspace{1.5ex}

\noindent
Lampiran ini memuat seluruh tautan ke berkas pendukung penelitian: kode sumber,
data, tutorial menjalankan, panduan pengguna, dan bahan presentasi. Seluruh
berkas yang dapat dipublikasikan berada pada satu repositori GitHub; berkas
yang berukuran besar atau tidak cocok disimpan di repositori kode (rekaman
demo, kumpulan slide) berada pada Google Drive dan tautannya dicantumkan di
Lampiran~B.

% ---------------------------------------------------------------------------
\lampiransec{A}{Repositori Kode Sumber dan Berkas Pendukung}

\noindent
\textbf{Repositori utama:} \url{\tautanRepo}

\vspace{1ex}
\noindent
Seluruh butir pada tabel berikut berada di dalam repositori tersebut. Kolom
``Lokasi'' adalah lintasan relatif terhadap akar repositori.

\vspace{1ex}
\renewcommand{\arraystretch}{1.25}
\begin{longtable}{@{}p{0.28\textwidth}p{0.27\textwidth}p{0.41\textwidth}@{}}
\toprule
\textbf{Butir lampiran} & \textbf{Lokasi di repositori} & \textbf{Keterangan} \\
\midrule
\endfirsthead
\multicolumn{3}{@{}l}{\footnotesize\itshape Lampiran A, lanjutan}\\
\toprule
\textbf{Butir lampiran} & \textbf{Lokasi di repositori} & \textbf{Keterangan} \\
\midrule
\endhead
\midrule
\multicolumn{3}{r@{}}{\footnotesize\itshape bersambung ke halaman berikutnya}\\
\endfoot
\bottomrule
\endlastfoot

\textbf{Kode sumber modul OpenMRS} &
\texttt{backend/} \newline
\texttt{~~openmrs-module-} \newline
\texttt{~~tfidf-search/} &
Modul Java (\texttt{.omod}). Berisi seluruh algoritma K1--K6: normalisasi,
tokenisasi kata dan kepingan karakter, pembobotan \emph{ltc}, fusi tingkat
satu, \emph{Weighted RRF}, baseline B0, jalur saran ketik, serta dua endpoint
REST. Dilengkapi 47 uji unit. \\

\textbf{Kode sumber antarmuka} &
\texttt{frontend/} \newline
\texttt{~~esm-unified-search/} &
Microfrontend (ESM) untuk OpenMRS Reference Application~3: halaman
``Pencarian Terpadu'', halaman ``Perbandingan Pencarian'', dan panel
evaluasi. \\

\textbf{Kode eksperimen} &
\texttt{riset/} &
Implementasi rujukan Python untuk seluruh angka pada naskah:
\texttt{eksperimen2.py} (sistem B0--E4, himpunan uji),
\texttt{eksperimen3\_baseline\_asli.py} (baseline B0$'$),
\texttt{eksperimen\_k2.py} (jalur saran ketik),
\texttt{sapuan\_dev.py} dan \texttt{sapuan\_alpha\_dev.py} (sapuan parameter). \\

\textbf{Dataset / basis data} &
\texttt{riset/data/*.jsonl} \newline
\texttt{riset/ekspor\_*.sql} &
Korpus hasil ekspor dari basis data OpenMRS dalam format JSON Lines
(\texttt{konsep}, \texttt{obat}, \texttt{pasien}, \texttt{lain},
\texttt{hasillab}, \texttt{kondisi}), beserta kueri SQL yang dipakai untuk
mengekspornya sehingga dapat dibangkitkan ulang dari instalasi OpenMRS mana
pun. \\

\textbf{Berkas hasil eksperimen} &
\texttt{riset/hasil4/} \dots \texttt{riset/hasil6/} &
Keluaran mentah tiap eksperimen (\texttt{hasil.json},
\texttt{per\_query.json}, \texttt{ringkasan.csv}, \texttt{laporan.md}).
Pemetaan berkas ke tabel dan gambar naskah ada pada Lampiran~F. \\

\textbf{Tutorial menjalankan kode} &
\texttt{README.md} &
Panduan lengkap dari nol: prasyarat, penyalaan stack Docker, pembangunan
modul, pemasangan, verifikasi, sampai halaman terbuka di peramban. Disertai
bagian penanganan galat yang sering muncul. \\

\textbf{Panduan reproduksi hasil} &
\texttt{docs/reproduksi.md} &
Sembilan langkah untuk mengulang seluruh angka penelitian, masing-masing
dengan kriteria ``cara tahu berhasil'' yang dapat diperiksa. \\

\textbf{Panduan pengguna} &
\texttt{docs/panduan-pengguna.md} &
Cara memakai halaman Pencarian Terpadu, halaman Perbandingan Pencarian,
kotak saran ketik, dan panel evaluasi dari sudut pandang pengguna akhir. \\

\textbf{Spesifikasi algoritma} &
\texttt{docs/algoritma.md} &
Spesifikasi K1--K6, baseline B0, dan jalur saran ketik K2, cukup rinci untuk
diimplementasikan ulang tanpa menebak. \\

\textbf{Kontrak data} &
\texttt{docs/kontrak-data.md} &
Bentuk dokumen virtual yang menyatukan delapan entitas OpenMRS; wajib diikuti
persis bila menambahkan entitas baru. \\

\textbf{Catatan keputusan} &
\texttt{docs/keputusan.md} &
Seluruh keputusan rancangan dan temuan penelitian secara kronologis, termasuk
kekeliruan yang ditemukan dan cara memperbaikinya. Rujukan utama untuk
pengembangan lanjutan. \\

\textbf{Fakta lingkungan} &
\texttt{docs/lingkungan.md} &
Versi platform terverifikasi, alamat, kredensial basis data, jumlah baris
tabel, dan jebakan lingkungan yang sudah terdokumentasi. \\

\textbf{Aturan proyek} &
\texttt{CLAUDE.md} &
Sepuluh aturan yang mengikat seluruh perubahan kode dan dokumen, termasuk
aturan determinisme dan larangan mengubah angka hasil. \\

\textbf{Proposal penelitian} &
\texttt{docs/proposal.html} &
Proposal lengkap beserta \emph{mockup} interaktif dan contoh perhitungan. \\

\textbf{Naskah artikel ini} &
\texttt{article/} &
Sumber \LaTeX{} naskah, skrip pembangun gambar, dan berkas hasil kompilasi. \\

\end{longtable}

% ---------------------------------------------------------------------------
\vspace{3ex}
\lampiransec{B}{Tautan Berkas Pendukung Eksternal}

\noindent
Berkas berikut berada di Google Drive karena ukurannya tidak sesuai untuk
disimpan pada repositori kode.

\vspace{1ex}
\renewcommand{\arraystretch}{1.6}
\begin{tabular}{@{}p{0.24\textwidth}p{0.72\textwidth}@{}}
\toprule
\textbf{Berkas} & \textbf{Tautan} \\
\midrule

\textbf{Kumpulan PPT Progress} \newline
{\footnotesize Seluruh slide laporan kemajuan penelitian, berurutan menurut
tanggal.} &
\tautanPPT \\

\textbf{Demo Aplikasi} \newline
{\footnotesize Rekaman layar modul yang berjalan: kueri gagal pada kotak
pencarian bawaan, kueri yang sama berhasil pada sistem usulan.} &
\tautanDemo \\

\bottomrule
\end{tabular}

\vspace{1.5ex}
\noindent
{\footnotesize \textbf{Catatan bagi penulis:} kedua tautan di atas diisi pada
blok \texttt{TAUTAN} di kepala berkas \texttt{article/lampiran.tex}
(makro \texttt{\textbackslash tautanPPT} dan \texttt{\textbackslash
tautanDemo}). Tidak perlu mengubah tabel ini.}

% ---------------------------------------------------------------------------
\vspace{3ex}
\lampiransec{C}{Peta Isi Repositori}

\begin{verbatim}
openmrs-tfidf-ngram/
|
+-- README.md                 tutorial lengkap menjalankan dari nol
+-- CLAUDE.md                 aturan proyek yang mengikat
|
+-- backend/openmrs-module-tfidf-search/
|   +-- api/                  algoritma inti + 47 uji unit
|   |   +-- src/main/java/org/openmrs/module/unifiedsearch/
|   |   |   +-- TextNormalizer.java        normalisasi teks
|   |   |   +-- Tokenizer.java             token kata + kepingan karakter
|   |   |   +-- TfIdfIndex.java            indeks terbalik, pembobotan ltc
|   |   |   +-- FusionSearch.java          fusi tingkat 1 (alpha)
|   |   |   +-- WeightedRrf.java           fusi tingkat 2 (Weighted RRF)
|   |   |   +-- OpenMrsHeuristic.java      baseline B0
|   |   |   +-- BigramJaccardSuggester.java  jalur saran ketik
|   |   |   +-- RankingEngine.java         pemilih mode b0/b1/e1/e3
|   |   |   +-- IndexBuilder.java          pembangunan 13 indeks saat startup
|   |   |   +-- EvalService.java           endpoint evaluasi
|   |   |   +-- source/                    proyeksi 8 tabel OpenMRS
|   |   +-- src/main/resources/gold-dev-100.json   standar emas 100 kueri dev
|   +-- omod/                 lapisan web: REST controller, halaman JSP
|
+-- frontend/esm-unified-search/   microfrontend RefApp 3
|
+-- riset/                    eksperimen Python
|   +-- eksperimen2.py                sistem B0-E4, himpunan uji 180 kueri
|   +-- eksperimen3_baseline_asli.py  baseline B0' (OpenMRS asli)
|   +-- eksperimen_k2.py              jalur saran ketik
|   +-- sapuan_dev.py                 sapuan NGRAM / K_RRF / EPS
|   +-- sapuan_alpha_dev.py           sapuan ALPHA
|   +-- data/*.jsonl                  korpus hasil ekspor
|   +-- ekspor_*.sql                  kueri SQL pembangkit korpus
|   +-- hasil4..hasil6/               keluaran mentah tiap eksperimen
|
+-- docs/                     dokumentasi
|   +-- reproduksi.md         9 langkah mengulang hasil
|   +-- panduan-pengguna.md   panduan pengguna akhir
|   +-- algoritma.md          spesifikasi K1-K6 + B0 + K2
|   +-- kontrak-data.md       bentuk dokumen virtual
|   +-- keputusan.md          catatan keputusan dan temuan, kronologis
|   +-- lingkungan.md         versi, alamat, kredensial, jebakan
|   +-- proposal.html         proposal + mockup interaktif
|
+-- scripts/                  otomasi pemasangan (PowerShell)
|   +-- pasang-modul.ps1      build .omod -> salin -> restart -> tunggu sehat
|   +-- pasang-esm.ps1        build ESM -> sisipkan ke importmap
|   +-- verify-openmrs.ps1    pemeriksaan kesehatan menyeluruh
|
+-- article/                  naskah artikel ini
|   +-- main.tex, lampiran.tex, referensi.bib
|   +-- gambar/buat_gambar.py  membangun gambar dari berkas hasil
|
+-- tugas/                    14 berkas tugas berurutan, masing-masing
                              dengan kriteria uji "Selesai kalau"
\end{verbatim}

% ---------------------------------------------------------------------------
\newpage
\lampiransec{D}{Ringkasan Langkah Menjalankan}

\noindent
Ringkasan ini memuat perintah intinya saja. Penjelasan lengkap setiap langkah,
termasuk cara memastikan langkah itu berhasil dan penanganan galat yang sering
muncul, ada pada \texttt{README.md} di akar repositori.

\vspace{1ex}
\noindent\textbf{Prasyarat:} Docker Desktop, Git, JDK~17, Node.js~20+,
Python~3.9+, dan PowerShell (Windows). Rincian versi terverifikasi ada pada
\texttt{docs/lingkungan.md}.

\vspace{1.5ex}
\renewcommand{\arraystretch}{1.2}
\begin{tabular}{@{}p{0.05\textwidth}p{0.40\textwidth}p{0.51\textwidth}@{}}
\toprule
\textbf{No.} & \textbf{Langkah} & \textbf{Perintah} \\
\midrule
1 & Klon repositori penelitian &
\texttt{git clone https://github.com/} \newline
\texttt{~~stephenprasetyachrismawan/} \newline
\texttt{~~openmrs-tfidf-ngram.git} \\
2 & Klon stack Docker OpenMRS sejajar dengan repositori penelitian &
\texttt{git clone https://github.com/} \newline
\texttt{openmrs/openmrs-distro-} \newline
\texttt{referenceapplication.git} \\
3 & Nyalakan stack (tanpa \texttt{-p}) &
\texttt{cd openmrs-distro-referenceapplication} \newline
\texttt{docker compose up -d} \\
4 & Tunggu empat kontainer sehat &
\texttt{docker ps} \\
5 & Bangun dan pasang modul backend &
\texttt{.\textbackslash scripts\textbackslash pasang-modul.ps1} \\
6 & Bangun dan pasang antarmuka ESM &
\texttt{.\textbackslash scripts\textbackslash pasang-esm.ps1} \\
7 & Verifikasi menyeluruh &
\texttt{.\textbackslash scripts\textbackslash verify-openmrs.ps1} \\
8 & Buka di peramban &
\texttt{http://127.0.0.1/openmrs/spa/} \newline
\texttt{unified-search} \\
\bottomrule
\end{tabular}

\vspace{2ex}
\noindent
\textbf{Peringatan lingkungan.} Seluruh alamat memakai \texttt{http://127.0.0.1},
\emph{bukan} \texttt{localhost}. Pada sebagian mesin Windows, \texttt{localhost}
teresolusi ke IPv6 \texttt{::1} yang mungkin dikuasai server lain, sehingga
seluruh lintasan \texttt{/openmrs/\dots} mengembalikan 404. Bila muncul halaman
bawaan Apache, yang salah adalah alamatnya, bukan modulnya.

\vspace{1ex}
\noindent
\textbf{Kredensial demo bawaan:} \texttt{admin} / \texttt{Admin123}.

% ---------------------------------------------------------------------------
\vspace{3ex}
\lampiransec{E}{Antarmuka REST Modul}

\noindent
Basis alamat: \texttt{http://127.0.0.1/openmrs/ws/rest/v1/unifiedsearch}

\vspace{1.5ex}
\renewcommand{\arraystretch}{1.3}
\begin{tabular}{@{}p{0.20\textwidth}p{0.34\textwidth}p{0.42\textwidth}@{}}
\toprule
\textbf{Endpoint} & \textbf{Parameter} & \textbf{Keterangan} \\
\midrule

\texttt{GET /unifiedsearch} &
\texttt{q} (wajib) \newline
\texttt{mode} = \texttt{b0} $|$ \texttt{b1} $|$ \texttt{e1} $|$ \texttt{e3}
  (bawaan \texttt{e3}) \newline
\texttt{limit} (bawaan 10) \newline
\texttt{entitas} (opsional) &
Pencarian terpadu. Parameter \texttt{mode} memilih sistem, sehingga kontribusi
tiap komponen dapat diukur ulang kapan saja tanpa mengubah kode. \\

\texttt{GET /unifiedsearch/} \newline \texttt{saran} &
\texttt{q} (wajib) \newline
\texttt{limit} (bawaan 6) &
Saran ketik ``Maksud Anda'' berbasis Jaccard bigram. Jalur terpisah; tidak
menyentuh peringkat. Hasil berjenis pasien, hasil laboratorium, dan kondisi
disaring menurut privilese \texttt{View Patients} pemanggil. \\

\texttt{GET /unifiedsearch/} \newline \texttt{eval} &
\texttt{mode} (bawaan \texttt{e3}) &
Menjalankan evaluasi atas 100 kueri pengembangan dan mengembalikan P@1, P@5,
R@10, MRR, MAP, nDCG@10, serta proporsi kueri tanpa hasil. Jawaban menyertakan
\texttt{gold\_sha256} agar standar emas yang dipakai dapat dicocokkan. \\

\bottomrule
\end{tabular}

\vspace{1.5ex}
\noindent
Setiap jawaban menyertakan tajuk \texttt{X-Unifiedsearch-Waktu-Ms} berisi waktu
pemrosesan. Waktu sengaja ditempatkan pada tajuk, bukan pada badan JSON, agar
dua panggilan yang sama menghasilkan badan jawaban yang identik pada tingkat
bita --- syarat pengujian determinisme.

\vspace{1.5ex}
\noindent\textbf{Contoh pemanggilan:}
\begin{verbatim}
curl -u admin:Admin123 \
  "http://127.0.0.1/openmrs/ws/rest/v1/unifiedsearch?q=diabete%20melitus&mode=e3"
\end{verbatim}

% ---------------------------------------------------------------------------
\vspace{2ex}
\lampiransec{F}{Pemetaan Berkas Hasil ke Tabel dan Gambar Naskah}

\noindent
Tidak ada angka pada naskah ini yang diketik ulang secara manual ke dalam
gambar. Skrip \texttt{article/gambar/buat\_gambar.py} membaca berkas hasil
di bawah dan membangun ulang seluruh gambar; bila eksperimen dijalankan ulang
dan angkanya berubah, gambarnya ikut berubah.

\vspace{1.5ex}
\renewcommand{\arraystretch}{1.25}
\begin{tabular}{@{}p{0.26\textwidth}p{0.30\textwidth}p{0.40\textwidth}@{}}
\toprule
\textbf{Bagian naskah} & \textbf{Berkas sumber} & \textbf{Cakupan} \\
\midrule
Tabel~\ref{tab:utama}, \ref{tab:ablasi}, \ref{tab:tipe}, \ref{tab:entitas};
Gambar~\ref{fig:sistem}, \ref{fig:tipe} &
\texttt{riset/hasil6/hasil.json} &
260 kueri uji, korpus 8 entitas, sekali jalan, parameter terkunci \\

Gambar~\ref{fig:sapuan} &
\texttt{riset/hasil6/sapuan\_dev.json} &
Sapuan $\alpha$, $n$, $k$, $\varepsilon$ pada 100 kueri pengembangan \\

Tabel~\ref{tab:b0prime}; Gambar~\ref{fig:b0prime} &
\texttt{riset/hasil4/hasil.json} &
34 kueri konsep, baseline B0$'$ (endpoint OpenMRS asli) \\

Tabel~\ref{tab:k2}; Gambar~\ref{fig:k2} &
\texttt{riset/hasil5/hasil.json} &
214 kueri, jalur saran ketik, korpus 8 entitas \\

Tabel~\ref{tab:verifikasi}, angka latensi dan determinisme &
\texttt{docs/keputusan.md} bagian C1--C4, F1--F2 &
Verifikasi silang Java--Python dan pengukuran kinerja modul \\
\bottomrule
\end{tabular}

\vspace{2ex}
\noindent
\textbf{Cara membangkitkan ulang seluruh angka:} jalankan
\texttt{python riset/eksperimen2.py} (himpunan uji),
\texttt{python riset/eksperimen3\_baseline\_asli.py} (baseline B0$'$, perlu
stack Docker berjalan), dan \texttt{python riset/eksperimen\_k2.py}
(jalur saran ketik), lalu \texttt{python article/gambar/buat\_gambar.py} untuk
membangun ulang gambar. Seluruh skrip deterministik: benih acak dipatok,
pengurutan memakai kunci majemuk, dan iterasi himpunan selalu diurutkan lebih
dahulu, sehingga dua proses terpisah menghasilkan berkas yang identik pada
tingkat bita.

\twocolumn

%%  tepat sebelum \end{document}, sesudah daftar pustaka.
%%
%%  >>> YANG PERLU ANDA ISI ADA DI BLOK "TAUTAN" DI BAWAH INI. <<<
%%      Ganti isi tiga makro berikut, lalu kompilasi ulang. Tidak ada tempat
%%      lain yang perlu disentuh -- seluruh tabel di bawah membaca makro ini.
%% ---------------------------------------------------------------------------

% ===== TAUTAN — ISI DI SINI =================================================

% Repositori GitHub (sudah terisi; ubah kalau repo dipindahkan).
\newcommand{\tautanRepo}{https://github.com/stephenprasetyachrismawan/openmrs-tfidf-ngram}

% Google Drive — kumpulan PPT progress. Ganti teks di dalam kurung kurawal.
\newcommand{\tautanPPT}{\textit{[tempelkan tautan Google Drive di sini]}}

% Google Drive — demo aplikasi (rekaman layar / berkas demo).
\newcommand{\tautanDemo}{\textit{[tempelkan tautan Google Drive di sini]}}

% ===== AKHIR BLOK TAUTAN ====================================================



% Judul bagian lampiran. Tidak memakai \section* bawaan ieeecolor.cls
% karena spasi vertikalnya dirancang untuk tata letak dua kolom dan
% meninggalkan lubang besar pada halaman satu kolom.
\newcommand{\lampiransec}[2]{%
  \par\addvspace{3ex}%
  \noindent{\large\bfseries Lampiran #1\quad #2}%
  \par\addvspace{1.2ex}%
  \noindent\ignorespaces}

\newpage
\onecolumn
\raggedbottom   % kelas IEEE memakai \flushbottom;
                % pada halaman lampiran yang isinya sedikit, itu meregangkan
                % jarak antar-paragraf sampai memenuhi tinggi halaman.

\begin{center}
{\Large\bfseries LAMPIRAN}\\[2pt]
{\normalsize Pencarian Terpadu Tahan Salah Ketik pada OpenMRS}
\end{center}

\vspace{1.5ex}

\noindent
Lampiran ini memuat seluruh tautan ke berkas pendukung penelitian: kode sumber,
data, tutorial menjalankan, panduan pengguna, dan bahan presentasi. Seluruh
berkas yang dapat dipublikasikan berada pada satu repositori GitHub; berkas
yang berukuran besar atau tidak cocok disimpan di repositori kode (rekaman
demo, kumpulan slide) berada pada Google Drive dan tautannya dicantumkan di
Lampiran~B.

% ---------------------------------------------------------------------------
\lampiransec{A}{Repositori Kode Sumber dan Berkas Pendukung}

\noindent
\textbf{Repositori utama:} \url{\tautanRepo}

\vspace{1ex}
\noindent
Seluruh butir pada tabel berikut berada di dalam repositori tersebut. Kolom
``Lokasi'' adalah lintasan relatif terhadap akar repositori.

\vspace{1ex}
\renewcommand{\arraystretch}{1.25}
\begin{longtable}{@{}p{0.28\textwidth}p{0.27\textwidth}p{0.41\textwidth}@{}}
\toprule
\textbf{Butir lampiran} & \textbf{Lokasi di repositori} & \textbf{Keterangan} \\
\midrule
\endfirsthead
\multicolumn{3}{@{}l}{\footnotesize\itshape Lampiran A, lanjutan}\\
\toprule
\textbf{Butir lampiran} & \textbf{Lokasi di repositori} & \textbf{Keterangan} \\
\midrule
\endhead
\midrule
\multicolumn{3}{r@{}}{\footnotesize\itshape bersambung ke halaman berikutnya}\\
\endfoot
\bottomrule
\endlastfoot

\textbf{Kode sumber modul OpenMRS} &
\texttt{backend/} \newline
\texttt{~~openmrs-module-} \newline
\texttt{~~tfidf-search/} &
Modul Java (\texttt{.omod}). Berisi seluruh algoritma K1--K6: normalisasi,
tokenisasi kata dan kepingan karakter, pembobotan \emph{ltc}, fusi tingkat
satu, \emph{Weighted RRF}, baseline B0, jalur saran ketik, serta dua endpoint
REST. Dilengkapi 47 uji unit. \\

\textbf{Kode sumber antarmuka} &
\texttt{frontend/} \newline
\texttt{~~esm-unified-search/} &
Microfrontend (ESM) untuk OpenMRS Reference Application~3: halaman
``Pencarian Terpadu'', halaman ``Perbandingan Pencarian'', dan panel
evaluasi. \\

\textbf{Kode eksperimen} &
\texttt{riset/} &
Implementasi rujukan Python untuk seluruh angka pada naskah:
\texttt{eksperimen2.py} (sistem B0--E4, himpunan uji),
\texttt{eksperimen3\_baseline\_asli.py} (baseline B0$'$),
\texttt{eksperimen\_k2.py} (jalur saran ketik),
\texttt{sapuan\_dev.py} dan \texttt{sapuan\_alpha\_dev.py} (sapuan parameter). \\

\textbf{Dataset / basis data} &
\texttt{riset/data/*.jsonl} \newline
\texttt{riset/ekspor\_*.sql} &
Korpus hasil ekspor dari basis data OpenMRS dalam format JSON Lines
(\texttt{konsep}, \texttt{obat}, \texttt{pasien}, \texttt{lain},
\texttt{hasillab}, \texttt{kondisi}), beserta kueri SQL yang dipakai untuk
mengekspornya sehingga dapat dibangkitkan ulang dari instalasi OpenMRS mana
pun. \\

\textbf{Berkas hasil eksperimen} &
\texttt{riset/hasil4/} \dots \texttt{riset/hasil6/} &
Keluaran mentah tiap eksperimen (\texttt{hasil.json},
\texttt{per\_query.json}, \texttt{ringkasan.csv}, \texttt{laporan.md}).
Pemetaan berkas ke tabel dan gambar naskah ada pada Lampiran~F. \\

\textbf{Tutorial menjalankan kode} &
\texttt{README.md} &
Panduan lengkap dari nol: prasyarat, penyalaan stack Docker, pembangunan
modul, pemasangan, verifikasi, sampai halaman terbuka di peramban. Disertai
bagian penanganan galat yang sering muncul. \\

\textbf{Panduan reproduksi hasil} &
\texttt{docs/reproduksi.md} &
Sembilan langkah untuk mengulang seluruh angka penelitian, masing-masing
dengan kriteria ``cara tahu berhasil'' yang dapat diperiksa. \\

\textbf{Panduan pengguna} &
\texttt{docs/panduan-pengguna.md} &
Cara memakai halaman Pencarian Terpadu, halaman Perbandingan Pencarian,
kotak saran ketik, dan panel evaluasi dari sudut pandang pengguna akhir. \\

\textbf{Spesifikasi algoritma} &
\texttt{docs/algoritma.md} &
Spesifikasi K1--K6, baseline B0, dan jalur saran ketik K2, cukup rinci untuk
diimplementasikan ulang tanpa menebak. \\

\textbf{Kontrak data} &
\texttt{docs/kontrak-data.md} &
Bentuk dokumen virtual yang menyatukan delapan entitas OpenMRS; wajib diikuti
persis bila menambahkan entitas baru. \\

\textbf{Catatan keputusan} &
\texttt{docs/keputusan.md} &
Seluruh keputusan rancangan dan temuan penelitian secara kronologis, termasuk
kekeliruan yang ditemukan dan cara memperbaikinya. Rujukan utama untuk
pengembangan lanjutan. \\

\textbf{Fakta lingkungan} &
\texttt{docs/lingkungan.md} &
Versi platform terverifikasi, alamat, kredensial basis data, jumlah baris
tabel, dan jebakan lingkungan yang sudah terdokumentasi. \\

\textbf{Aturan proyek} &
\texttt{CLAUDE.md} &
Sepuluh aturan yang mengikat seluruh perubahan kode dan dokumen, termasuk
aturan determinisme dan larangan mengubah angka hasil. \\

\textbf{Proposal penelitian} &
\texttt{docs/proposal.html} &
Proposal lengkap beserta \emph{mockup} interaktif dan contoh perhitungan. \\

\textbf{Naskah artikel ini} &
\texttt{article/} &
Sumber \LaTeX{} naskah, skrip pembangun gambar, dan berkas hasil kompilasi. \\

\end{longtable}

% ---------------------------------------------------------------------------
\vspace{3ex}
\lampiransec{B}{Tautan Berkas Pendukung Eksternal}

\noindent
Berkas berikut berada di Google Drive karena ukurannya tidak sesuai untuk
disimpan pada repositori kode.

\vspace{1ex}
\renewcommand{\arraystretch}{1.6}
\begin{tabular}{@{}p{0.24\textwidth}p{0.72\textwidth}@{}}
\toprule
\textbf{Berkas} & \textbf{Tautan} \\
\midrule

\textbf{Kumpulan PPT Progress} \newline
{\footnotesize Seluruh slide laporan kemajuan penelitian, berurutan menurut
tanggal.} &
\tautanPPT \\

\textbf{Demo Aplikasi} \newline
{\footnotesize Rekaman layar modul yang berjalan: kueri gagal pada kotak
pencarian bawaan, kueri yang sama berhasil pada sistem usulan.} &
\tautanDemo \\

\bottomrule
\end{tabular}

\vspace{1.5ex}
\noindent
{\footnotesize \textbf{Catatan bagi penulis:} kedua tautan di atas diisi pada
blok \texttt{TAUTAN} di kepala berkas \texttt{article/lampiran.tex}
(makro \texttt{\textbackslash tautanPPT} dan \texttt{\textbackslash
tautanDemo}). Tidak perlu mengubah tabel ini.}

% ---------------------------------------------------------------------------
\vspace{3ex}
\lampiransec{C}{Peta Isi Repositori}

\begin{verbatim}
openmrs-tfidf-ngram/
|
+-- README.md                 tutorial lengkap menjalankan dari nol
+-- CLAUDE.md                 aturan proyek yang mengikat
|
+-- backend/openmrs-module-tfidf-search/
|   +-- api/                  algoritma inti + 47 uji unit
|   |   +-- src/main/java/org/openmrs/module/unifiedsearch/
|   |   |   +-- TextNormalizer.java        normalisasi teks
|   |   |   +-- Tokenizer.java             token kata + kepingan karakter
|   |   |   +-- TfIdfIndex.java            indeks terbalik, pembobotan ltc
|   |   |   +-- FusionSearch.java          fusi tingkat 1 (alpha)
|   |   |   +-- WeightedRrf.java           fusi tingkat 2 (Weighted RRF)
|   |   |   +-- OpenMrsHeuristic.java      baseline B0
|   |   |   +-- BigramJaccardSuggester.java  jalur saran ketik
|   |   |   +-- RankingEngine.java         pemilih mode b0/b1/e1/e3
|   |   |   +-- IndexBuilder.java          pembangunan 13 indeks saat startup
|   |   |   +-- EvalService.java           endpoint evaluasi
|   |   |   +-- source/                    proyeksi 8 tabel OpenMRS
|   |   +-- src/main/resources/gold-dev-100.json   standar emas 100 kueri dev
|   +-- omod/                 lapisan web: REST controller, halaman JSP
|
+-- frontend/esm-unified-search/   microfrontend RefApp 3
|
+-- riset/                    eksperimen Python
|   +-- eksperimen2.py                sistem B0-E4, himpunan uji 180 kueri
|   +-- eksperimen3_baseline_asli.py  baseline B0' (OpenMRS asli)
|   +-- eksperimen_k2.py              jalur saran ketik
|   +-- sapuan_dev.py                 sapuan NGRAM / K_RRF / EPS
|   +-- sapuan_alpha_dev.py           sapuan ALPHA
|   +-- data/*.jsonl                  korpus hasil ekspor
|   +-- ekspor_*.sql                  kueri SQL pembangkit korpus
|   +-- hasil4..hasil6/               keluaran mentah tiap eksperimen
|
+-- docs/                     dokumentasi
|   +-- reproduksi.md         9 langkah mengulang hasil
|   +-- panduan-pengguna.md   panduan pengguna akhir
|   +-- algoritma.md          spesifikasi K1-K6 + B0 + K2
|   +-- kontrak-data.md       bentuk dokumen virtual
|   +-- keputusan.md          catatan keputusan dan temuan, kronologis
|   +-- lingkungan.md         versi, alamat, kredensial, jebakan
|   +-- proposal.html         proposal + mockup interaktif
|
+-- scripts/                  otomasi pemasangan (PowerShell)
|   +-- pasang-modul.ps1      build .omod -> salin -> restart -> tunggu sehat
|   +-- pasang-esm.ps1        build ESM -> sisipkan ke importmap
|   +-- verify-openmrs.ps1    pemeriksaan kesehatan menyeluruh
|
+-- article/                  naskah artikel ini
|   +-- main.tex, lampiran.tex, referensi.bib
|   +-- gambar/buat_gambar.py  membangun gambar dari berkas hasil
|
+-- tugas/                    14 berkas tugas berurutan, masing-masing
                              dengan kriteria uji "Selesai kalau"
\end{verbatim}

% ---------------------------------------------------------------------------
\newpage
\lampiransec{D}{Ringkasan Langkah Menjalankan}

\noindent
Ringkasan ini memuat perintah intinya saja. Penjelasan lengkap setiap langkah,
termasuk cara memastikan langkah itu berhasil dan penanganan galat yang sering
muncul, ada pada \texttt{README.md} di akar repositori.

\vspace{1ex}
\noindent\textbf{Prasyarat:} Docker Desktop, Git, JDK~17, Node.js~20+,
Python~3.9+, dan PowerShell (Windows). Rincian versi terverifikasi ada pada
\texttt{docs/lingkungan.md}.

\vspace{1.5ex}
\renewcommand{\arraystretch}{1.2}
\begin{tabular}{@{}p{0.05\textwidth}p{0.40\textwidth}p{0.51\textwidth}@{}}
\toprule
\textbf{No.} & \textbf{Langkah} & \textbf{Perintah} \\
\midrule
1 & Klon repositori penelitian &
\texttt{git clone https://github.com/} \newline
\texttt{~~stephenprasetyachrismawan/} \newline
\texttt{~~openmrs-tfidf-ngram.git} \\
2 & Klon stack Docker OpenMRS sejajar dengan repositori penelitian &
\texttt{git clone https://github.com/} \newline
\texttt{openmrs/openmrs-distro-} \newline
\texttt{referenceapplication.git} \\
3 & Nyalakan stack (tanpa \texttt{-p}) &
\texttt{cd openmrs-distro-referenceapplication} \newline
\texttt{docker compose up -d} \\
4 & Tunggu empat kontainer sehat &
\texttt{docker ps} \\
5 & Bangun dan pasang modul backend &
\texttt{.\textbackslash scripts\textbackslash pasang-modul.ps1} \\
6 & Bangun dan pasang antarmuka ESM &
\texttt{.\textbackslash scripts\textbackslash pasang-esm.ps1} \\
7 & Verifikasi menyeluruh &
\texttt{.\textbackslash scripts\textbackslash verify-openmrs.ps1} \\
8 & Buka di peramban &
\texttt{http://127.0.0.1/openmrs/spa/} \newline
\texttt{unified-search} \\
\bottomrule
\end{tabular}

\vspace{2ex}
\noindent
\textbf{Peringatan lingkungan.} Seluruh alamat memakai \texttt{http://127.0.0.1},
\emph{bukan} \texttt{localhost}. Pada sebagian mesin Windows, \texttt{localhost}
teresolusi ke IPv6 \texttt{::1} yang mungkin dikuasai server lain, sehingga
seluruh lintasan \texttt{/openmrs/\dots} mengembalikan 404. Bila muncul halaman
bawaan Apache, yang salah adalah alamatnya, bukan modulnya.

\vspace{1ex}
\noindent
\textbf{Kredensial demo bawaan:} \texttt{admin} / \texttt{Admin123}.

% ---------------------------------------------------------------------------
\vspace{3ex}
\lampiransec{E}{Antarmuka REST Modul}

\noindent
Basis alamat: \texttt{http://127.0.0.1/openmrs/ws/rest/v1/unifiedsearch}

\vspace{1.5ex}
\renewcommand{\arraystretch}{1.3}
\begin{tabular}{@{}p{0.20\textwidth}p{0.34\textwidth}p{0.42\textwidth}@{}}
\toprule
\textbf{Endpoint} & \textbf{Parameter} & \textbf{Keterangan} \\
\midrule

\texttt{GET /unifiedsearch} &
\texttt{q} (wajib) \newline
\texttt{mode} = \texttt{b0} $|$ \texttt{b1} $|$ \texttt{e1} $|$ \texttt{e3}
  (bawaan \texttt{e3}) \newline
\texttt{limit} (bawaan 10) \newline
\texttt{entitas} (opsional) &
Pencarian terpadu. Parameter \texttt{mode} memilih sistem, sehingga kontribusi
tiap komponen dapat diukur ulang kapan saja tanpa mengubah kode. \\

\texttt{GET /unifiedsearch/} \newline \texttt{saran} &
\texttt{q} (wajib) \newline
\texttt{limit} (bawaan 6) &
Saran ketik ``Maksud Anda'' berbasis Jaccard bigram. Jalur terpisah; tidak
menyentuh peringkat. Hasil berjenis pasien, hasil laboratorium, dan kondisi
disaring menurut privilese \texttt{View Patients} pemanggil. \\

\texttt{GET /unifiedsearch/} \newline \texttt{eval} &
\texttt{mode} (bawaan \texttt{e3}) &
Menjalankan evaluasi atas 100 kueri pengembangan dan mengembalikan P@1, P@5,
R@10, MRR, MAP, nDCG@10, serta proporsi kueri tanpa hasil. Jawaban menyertakan
\texttt{gold\_sha256} agar standar emas yang dipakai dapat dicocokkan. \\

\bottomrule
\end{tabular}

\vspace{1.5ex}
\noindent
Setiap jawaban menyertakan tajuk \texttt{X-Unifiedsearch-Waktu-Ms} berisi waktu
pemrosesan. Waktu sengaja ditempatkan pada tajuk, bukan pada badan JSON, agar
dua panggilan yang sama menghasilkan badan jawaban yang identik pada tingkat
bita --- syarat pengujian determinisme.

\vspace{1.5ex}
\noindent\textbf{Contoh pemanggilan:}
\begin{verbatim}
curl -u admin:Admin123 \
  "http://127.0.0.1/openmrs/ws/rest/v1/unifiedsearch?q=diabete%20melitus&mode=e3"
\end{verbatim}

% ---------------------------------------------------------------------------
\vspace{2ex}
\lampiransec{F}{Pemetaan Berkas Hasil ke Tabel dan Gambar Naskah}

\noindent
Tidak ada angka pada naskah ini yang diketik ulang secara manual ke dalam
gambar. Skrip \texttt{article/gambar/buat\_gambar.py} membaca berkas hasil
di bawah dan membangun ulang seluruh gambar; bila eksperimen dijalankan ulang
dan angkanya berubah, gambarnya ikut berubah.

\vspace{1.5ex}
\renewcommand{\arraystretch}{1.25}
\begin{tabular}{@{}p{0.26\textwidth}p{0.30\textwidth}p{0.40\textwidth}@{}}
\toprule
\textbf{Bagian naskah} & \textbf{Berkas sumber} & \textbf{Cakupan} \\
\midrule
Tabel~\ref{tab:utama}, \ref{tab:ablasi}, \ref{tab:tipe}, \ref{tab:entitas};
Gambar~\ref{fig:sistem}, \ref{fig:tipe} &
\texttt{riset/hasil6/hasil.json} &
260 kueri uji, korpus 8 entitas, sekali jalan, parameter terkunci \\

Gambar~\ref{fig:sapuan} &
\texttt{riset/hasil6/sapuan\_dev.json} &
Sapuan $\alpha$, $n$, $k$, $\varepsilon$ pada 100 kueri pengembangan \\

Tabel~\ref{tab:b0prime}; Gambar~\ref{fig:b0prime} &
\texttt{riset/hasil4/hasil.json} &
34 kueri konsep, baseline B0$'$ (endpoint OpenMRS asli) \\

Tabel~\ref{tab:k2}; Gambar~\ref{fig:k2} &
\texttt{riset/hasil5/hasil.json} &
214 kueri, jalur saran ketik, korpus 8 entitas \\

Tabel~\ref{tab:verifikasi}, angka latensi dan determinisme &
\texttt{docs/keputusan.md} bagian C1--C4, F1--F2 &
Verifikasi silang Java--Python dan pengukuran kinerja modul \\
\bottomrule
\end{tabular}

\vspace{2ex}
\noindent
\textbf{Cara membangkitkan ulang seluruh angka:} jalankan
\texttt{python riset/eksperimen2.py} (himpunan uji),
\texttt{python riset/eksperimen3\_baseline\_asli.py} (baseline B0$'$, perlu
stack Docker berjalan), dan \texttt{python riset/eksperimen\_k2.py}
(jalur saran ketik), lalu \texttt{python article/gambar/buat\_gambar.py} untuk
membangun ulang gambar. Seluruh skrip deterministik: benih acak dipatok,
pengurutan memakai kunci majemuk, dan iterasi himpunan selalu diurutkan lebih
dahulu, sehingga dua proses terpisah menghasilkan berkas yang identik pada
tingkat bita.

\twocolumn
